\subsection{Mimo}

První analyzovanou aplikací je aplikace Mimo, která se nejvíce blíží vizi tohoto projektu a bude tak přímou konkurencí vznikající aplikace. Mimo se zaměřuje zejména na výuku moderních programovacích jazyků a frameworků. V době psaní této práce Mimo nabízí 9 kurzů, 4 tzv. kariérní cesty a má napříč platformami přes 40 milionů registrovaných uživatelů. Mimo je k dispozici jako webová aplikace a dále také jako mobilní aplikace na platformách Android a iOS. Pro analýzu použiji aplikaci na platformě Android.

% Write
\subsubsection{Vzhled a uživatelská zkušenost}

% Odstavec: Obecný dojem (komplexita/přehlednost UI, paleta barev, typografie)
Mimo se od dalších konkurentů liší zejména svým vzhledem uživatelského rozhraní a stylem výuky. Uživatelské rozhraní je velice minimalistické a přehledné. Místy může aplikace působit až příliš jednoduchým dojmem. Aplikace využívá fialovou monochromatickou paletu barev. V rámci typografie Mimo využívá neproporciální písmo % Write - vysvětlit buď tady nebo v pojmech
pro některé nadpisy. To lze vidět zejména v mobilní aplikaci a na webových stránkách. Pro textové prvky pak využívá bezpatkové písmo Aeonik.

% Odstavec: Rozložení prvků na obrazovce (lekce v kurzu, učení)
Rozložení prvků na obrazovce je přehledné. Na hlavní stránce kurzu jsou umístěny jednotlivé lekce, které jsou vyobrazeny ve formě tlačítek na čáře připomínající cestu kurzem. Až po kliknutí na danou lekci se zobrazí její název a možnost spuštění dané lekce. Toto vyobrazení lekcí může na uživatele působit vizuálně poutavým dojmem, nicméně může být obtížné se pomocí něj orientovat v kurzu, protože o lekcích nejsou na první pohled zobrazeny žádné podrobnější informace.

% Odstavec: Navigace v aplikaci
Pro navigaci v mobilní aplikaci používá Mimo velice jednoduchou spodní navigační lištu % Write - zde ověřit jestli se to tak jmenuje v češtině
s pouze pěti prvky. Díky tomu je navigace v aplikaci velice snadná a rychlá.

% Odstavec: Srovnání webového a mobilního prostředí
Webová verze aplikace Mimo je velice podobná mobilní verzi. Rozdíl je především v tom, že ve webové aplikaci jsou lekce vyobrazeny ve formě seznamu. Pro navigaci je použita horní navigační lišta a pro navigaci mezi kapitolami kurzu pak boční navigační lišta.

% Odstavec: Animace a efekty a mikro interakce
V aplikaci jsou použity animace a efekty velice zřídka. Některé prvky uživatelského rozhraní jsou animovány při jejich zobrazení. Při učení jsou také použity animace pro přechody mezi jednotlivými cvičeními. 

% Write
\subsubsection{Způsoby učení}

% Struktura kurzu (lineární, tematické moduly, career paths, kapitoly atd.)
Jednotlivé kurzy jsou uspořádány do sekcí, které obsahují jednotlivé lekce. Lekce jsou uspořádány lineárně a uživatel se je musí učit v daném pořadí. Kromě kurzů Mimo nabízí tzv. kariérní cesty, které mají podobnou strukturu jako kurzy, nicméně obsahují více témat.

% Jak se spustí učení a jak vypadá study session
Učení lze spustit kliknutím na danou lekci. Při učení se uživateli zobrazuje sekvence stránek na nichž se nachází buď vysvětlení dané látky nebo cvičení na procvičení dané látky. Při vyplňování cvičení je typicky uživateli zobrazena otázka, na kterou musí odpovědět například vyplněním textového pole v kódu. Po vyplnění odpovědi se uživateli zobrazí, zda byla odpověď správná a typicky se mu zobrazí výstup takového kódu, kdyby byl proveden. V kurzu SQL dotazů uživatel může například vidět tabulku s vybranými daty.

% Jak probíhá učení (jak vypadá study session, zpětná vazba okamžitá/opožděná (jak), vysvětlení chyb)
Při učení uživatel má k dispozici 5 srdíček, o která přichází, poku

% Druhy cvičení
V aplikaci 


% Write
\subsubsection{Způsoby motivace uživatele}

% Jaké formy gamifikace aplikace používá
% Jestli používá remindery a notifikace na učení


% Write
\subsubsection{Způsoby monetizace}

% Jaký model využívá
% Jaké všechny výhody mají premium uživatelé

% Write
\subsubsection{Silné a slabé stránky}

% Silné stránky
% Slabé stránky

% Review
Možnou nevýhodou minimalistického vzhledu uživatelského rozhraní může být pocit, který aplikace svým minimalismem uživatelům vyvolává. Aplikace místy může působit až příliš jednoduše a vyvolat tak dojem, že je příliš jednoduchá na to, aby vyučovala komplexní témata, kterým je například softwarové inženýrství.  

% Write
\subsubsection{Další poznatky}

% Zajímavé funkce
% (Zajímavé detaily)
